\begin{figure}[htbp]
	\centering
	\begin{tikzpicture}[>=Stealth, line width=0.8pt, scale=0.80]
		
		% Coordinate principali
		\coordinate (A)   at (0, 0);
		\coordinate (Sm)  at (2.5, 0);
		\coordinate (Sp)  at (5.5, 0);
		\coordinate (B)   at (8, 0);
		
		% Linee d'asse dei due tratti
		\draw[thick] (A) -- (Sm);
		\draw[thick] (Sp) -- (B);
		
		% Tratteggio fibra inferiore
		\draw[dashed, thick] (0,-0.2) -- (2.5,-0.2);
		\draw[dashed, thick] (5.5,-0.2) -- (8,-0.2);
		
		% Etichette estremi e sezione
		\node[above=4pt] at (A)  {$A$};
		\node[above=4pt] at (B)  {$B$};
		\node[above=6pt] at (4.0, 0) {$S$};
		
		% Linea di taglio in S (verticale tratteggiata)
		\draw[dashed, gray!60, thin] (4.0,-1.4) -- (4.0, 1.6);
		
		% Etichette tratti
		\node at (1.2, 0.5) {$\mathcal{C}^-$};
		\node at (6.8, 0.5) {$\mathcal{C}^+$};
		
		% Facce di taglio (linee verticali)
		\draw[thick] (2.5,-1.4) -- (2.5, 0.8);
		\draw[thick] (5.5,-0.8) -- (5.5, 1.4);
		
		% ── Faccia destra di C^- ─────────────────────────
		% N: verso destra
		\draw[->, line width=1.6pt]	(2.5, 0) -- (3.7, 0);
		\node[above=2pt] at (3.1, 0) {$N$};
		
		% T: verso il basso
		\draw[->, line width=1.6pt]	(2.5, 0) -- (2.5, -1.2);
		\node[right=2pt] at (2.5, -0.6) {$T$};
		
		% M: quarto di cerchio da T a N
		\draw[->, line width=1.6pt]	(2.5, -1.2) arc[start angle=270, end angle=360, radius=1.2];
		\node at (3.5, -1.1) {$M$};
		
		% ── Faccia sinistra di C^+ ────────────────────────
		% N: verso sinistra
		\draw[->, line width=1.6pt]	(5.5, 0) -- (4.3, 0);
		\node[above=2pt] at (4.9, 0) {$N$};
		
		% T: verso l'alto
		\draw[->, line width=1.6pt]
		(5.5, 0) -- (5.5, 1.2);
		\node[right=2pt] at (5.5, 0.6) {$T$};
		
		% M: quarto di cerchio da N a T
		\draw[->, line width=1.6pt]	(4.3, 0) arc[start angle=180,	end angle=90, radius=1.2];
		\node at (4.5, 1.1) {$M$};
		
		% Punti agli estremi
		\fill (A) circle (2pt);
		\fill (B) circle (2pt);
		
	\end{tikzpicture}
	\caption{Caratteristiche della sollecitazione
		nella sezione $S$: sulla faccia destra di
		$\mathcal{C}^-$ (normale uscente $+z^e$)
		agiscono $N$, $T$, $M$ nel verso positivo
		convenuto; sulla faccia sinistra di
		$\mathcal{C}^+$ le stesse grandezze agiscono
		con verso opposto per il terzo principio.}
	\label{fig:NTM_sezione}
\end{figure}
